
\medskip
\noindent
\SGARefTarget{sga7.sga7i.eix.n2-5}{2.5}. The orthogonality theorem has an equivalent formulation using a
polarization $\xi$ of $A_K$ and the alternating bilinear form (\SGARefLink{sga7.sga7i.eix.n1-0-8}{1.0.8}),
compatible with the operations of $\pi$,
\begin{equation}\tag{\SGARefTarget{sga7.sga7i.eix.n2-5-1}{2.5.1}}
 \varphi_\xi:T_\ell(A(\overline K))\times T_\ell(A(\overline K))
 \longrightarrow T_\ell(\overline K^*),
\end{equation}
obtained from the form $\varphi$ of (\SGARefLink{sga7.sga7i.eix.n1-0-3}{1.0.3}) by transporting structure in
the second argument along the isogeny of $\mathbb Z_\ell$-modules
\[
 T_\ell(A(\overline K))\longrightarrow T_\ell(A'(\overline K))
\]
induced by the isogeny $\psi_\xi:A_K\longrightarrow A'_K$ defined by
$\xi$.  The equality in (\SGARefLink{sga7.sga7i.eix.n2-4-6}{2.4.6}) then takes the simple form
\begin{equation}\tag{\SGARefTarget{sga7.sga7i.eix.n2-5-2}{2.5.2}}
 W=V\cap V^\perp,
\end{equation}
where we have set
\begin{equation}\tag{\SGARefTarget{sga7.sga7i.eix.n2-5-3}{2.5.3}}
 \left\{
 \begin{array}{l}
 U=T_\ell(A(\overline K)),\qquad
 V=U^I=\text{the ``fixed part'' of }T_\ell(A(\overline K)),\\[2mm]
 \hspace{4em}W=\text{the toric part of }T_\ell(A(\overline K)),
 \end{array}
 \right.
\end{equation}
and $\perp$ denotes the orthogonal with respect to the polarization form
(\SGARefLink{sga7.sga7i.eix.n2-5-1}{2.5.1}).  These modules give the following diagram of inclusions:
\begin{equation}\tag{\SGARefTarget{sga7.sga7i.eix.n2-5-4}{2.5.4}}
\begin{tikzcd}[column sep=large,row sep=large]
 && V\arrow[dl,hook',"2\lambda"'] && \\
 U & W^\perp=V+V^\perp\arrow[l,hook',"\mu"']
 && W=V\cap V^\perp
    \arrow[ul,hook',"2\alpha"']
    \arrow[dl,hook,"2\lambda"]
 & 0\arrow[l,hook',"\mu"'] \\
 && V^\perp\arrow[ul,hook,"2\alpha"] &&
\end{tikzcd}
\qquad ,
\end{equation}
where the integers $2\alpha$, $2\lambda$, and $\mu$ are the ranks of
the quotients of consecutive modules in the diagram, with the notation
of (\SGARefLink{sga7.sga7i.eix.n2-1-9}{2.1.9}).  (That $W$ has rank $\mu$ is the definition of $\mu$; that
$V/W$ has rank $2\alpha$ follows from the second relation in (\SGARefLink{sga7.sga7i.eix.n2-1-11}{2.1.11})
together with the exact sequence (\SGARefLink{sga7.sga7i.eix.n2-1-6}{2.1.6}).  Hence, by orthogonality,
$\operatorname{rank}(U/W^\perp)=\mu$ and
$\operatorname{rank}(W^\perp/V^\perp)=2\alpha$.  Finally, the formulas
\begin{equation}\tag{\SGARefTarget{sga7.sga7i.eix.n2-5-5}{2.5.5}}
 \operatorname{rank}\bigl(V/(V\cap V^\perp)\bigr)
 =\operatorname{rank}(V/W)=2\alpha
\end{equation}
and the dual formula follow from the preceding ones, in view of
(\SGARefLink{sga7.sga7i.eix.n2-1-10}{2.1.10}) and $\operatorname{rank}U=n$.)

\medskip
\noindent
\underline{Remarks} \SGARefTarget{sga7.sga7i.eix.n2-6}{2.6}.\quad a) Here we obtained the orthogonality
theorem \SGARefLink{sga7.sga7i.eix.n2-4}{2.4} as an easy consequence of the theory of biextensions
developed in the preceding sections.  The \underline{inclusion}
(\SGARefLink{sga7.sga7i.eix.n2-4-6}{2.4.6}) can be proved without using the theory of biextensions, by an
arithmetic argument in the style of the proof of the ``monodromy
theorem'' given in Expos\'e I.  One reduces to a situation in which $S$
has finite residue field, where the orthogonality result follows from a
simple consideration of \underline{weights} for the Frobenius operation
(using A.~Weil's result on the eigenvalues of that Frobenius).  There is
no need to invoke resolution of singularities here: one may state the
desired result for group schemes $A$ over schemes $S$ under sufficiently
general conditions.  Unfortunately, it does not appear possible to
obtain by this method equality in the preceding inclusion.  That
equality will apparently be used in an essential way in the proof of
the ``semistable reduction theorem'' in the next section.

\smallskip
\noindent
b) The orthogonality theorem \SGARefLink{sga7.sga7i.eix.n2-4}{2.4} allows the most important properties
of the N\'{e}ron model of an abelian variety $A_K$ to be expressed in
terms of the representation of the inertia group $I$ on the Tate module
$T_\ell(A(\overline K))$, where $\ell$ is a fixed prime different from
the residual characteristic.  (The case $\ell=p$ will be considered
separately later, in section 5.)  As a first statement of this kind,
though one that does not use the orthogonality theorem, recall the
good-reduction criterion \SGARefLink{sga7.sga7i.eix.n2-2-9}{2.2.9}.  We shall give many further examples
below, beginning in the next section (\SGARefLink{sga7.sga7i.eix.n3-5}{3.5}).
